\section{Limit Passage Under the Integral Sign}
\label{sec:scambio_integrale_limite}%

In example~\ref{ex:gamma_eulero}, we defined a function $\Gamma(x)$ via the
integral in $dt$ of a function that depends on the parameter $x$:
\[
  \Gamma(x) = \int_0^{+\infty} 
  e^{-t}t^{x-1}\, dt.
\] 
If the integrand function $f(x,t)=e^{-t}t^{x-1}$ is continuous and
differentiable, it is natural to ask whether the integral function $\Gamma$ is
also continuous and differentiable. If it is differentiable, it is natural to
ask how the derivative can be calculated.

In general, these questions can be easily answered if we are able to interchange
the limit and integral operators. Indeed, we observe that
\begin{align*}
  \lim_{x\to x_0} \Gamma(x) 
  &= \lim_{x\to x_0} \int_0^{+\infty} e^{-t}t^{x-1} \, dt 
  \stackrel{?}= \int_0^{+\infty} \lim_{x\to x_0} e^{-t}t^{x-1}\, dt\\
  &= \int_0^{+\infty} e^{-t}t^{x_0-1}\, dt
  = \Gamma(x_0)
\end{align*}
and 
\begin{align*}
  \Gamma'(x) 
  &= \frac{d}{dx} \int_0^{+\infty}e^{-t}t^{x-1}\, dx
  \stackrel{?}= \int_0^{+\infty} \frac{d}{dx}e^{-t}t^{x-1}\, dt \\
  &= \int_0^{+\infty} e^{-t}t^{x-1}\ln t\, dt.
\end{align*}
In both cases, the only step not formally justified is the one marked with a
question mark. The purpose of this chapter is to determine some conditions that
allow the interchange of the limit with the integral and the interchange of the
derivative with the integral.

\begin{example}%
The sequence $f_n(x)=x^n$ from example~\ref{ex:puntuale_non_uniforme} converges
pointwise to the function that is $0$ on the interval $[0,1)$.

In this case, the limit can be interchanged with the integral:
\begin{align*}
\lim_{n\to +\infty} \int_{-1}^1 f_n(x)\, dx 
&= \lim_{n\to+\infty} \Enclose{\frac{x^{n+1}}{n+1}}_{-1}^1
 = \lim_{n\to+\infty} \frac{1}{n+1} = 0 \\
\int_{-1}^1 \lim_{n\to+\infty} f_n(x)\, dx 
&= \int_{-1}^1 f(x)\, dx = 0.
\end{align*}
\end{example}

The previous example tells us that the interchange between the limit and the
integral can be done in certain cases. Unfortunately, pointwise convergence is
not a sufficient hypothesis, as seen in the following example.

\begin{example}
Consider the sequence of functions $f_n\colon[0,1]\to\RR$ defined by 
\[
 f_n(x) = \frac{nx}{1+n^2x^4}.  
\]
Clearly, for every $x\in \RR$:
\[
  \lim_{n\to+\infty} \frac{nx}{1+n^2x^4} = 0.
\]
Thus, the sequence $f_n$ converges pointwise to the identically zero function
$f=0$. However, the interchange of the limit with the integral in this case is
not possible since it results in
\begin{align*}
\lim_{n\to +\infty} \int_0^1 f_n(x)\, dx 
&= \lim_{n\to +\infty} \int_0^1 \frac{nx}{1+n^2x^4}\, dx \\
&=\lim_{n\to +\infty} \frac 1 2 \Enclose{\arctg(nx^2)}_0^1 
 = \frac 1 2 \lim_{n\to +\infty} \arctg n = \frac \pi 4 
\end{align*}
while obviously 
\[
 \int_0^1 \lim_{n\to +\infty} f_n(x)\, dx = \int_0^1 0\, dx = 0.
\]
\end{example}

The reason why in the previous example it is not possible to pass to the limit
under the integral sign is due to the fact that the function $f_n(x)$, although
tending to zero at every point $x\in \RR$, becomes very large near $x=0$ as $n$
increases. In fact, the graph of the function $f_n(x)$ is obtained from the
graph of the function $f(x) = \frac{x}{1+x^4}$ by compressing the variable $x$
by a factor of $n$ and expanding the variable $y$ by the same factor $n$. This
ensures that the area under the graph remains unchanged. The function $f_n$ has
a maximum point at $x=\frac{1}{\sqrt[4]3 \sqrt n}$ (this can be verified by
studying the sign of the derivative) and at this point, it assumes a value that
tends to infinity as $n\to +\infty$.

However, if there is uniform convergence and the integral is not improper, the
following holds.

\begin{theorem}[Limit Passage Under the Integral Sign]%
  \label{th:scambio_limite_integrale}%
  \mymark{***}%
    Let $f_n,f\colon [a,b] \to \RR$ be bounded and Riemann integrable functions
  on the closed and bounded interval $[a,b]$. Suppose further that there is
  uniform convergence: $f_n \rightrightarrows f$. Then
  \[
      \lim_{n\to+\infty} \int_a^b f_n(x)\, dx 
      = \int_a^b f(x)\, dx.
  \]

  Moreover, for any chosen $c\in [a,b]$, let
  \[
    F_k(x) = \int_{c}^x f_k(t)\, dt,
    \qquad
    F(x) = \int_{c}^x f(t)\, dt
  \]
  then there is uniform convergence $F_k \rightrightarrows F$.
\end{theorem}
  %
\begin{proof}
    If $f_n$ and $f$ are Riemann-integrable, then $f_n-f$ and $\abs{f_n-f}$ are
  also integrable (theorem~\ref{th:reticolo}) and we have:
  \begin{align*}
    \abs{\int_a^b f_n(x)\, dx - \int_a^b f(x)\, dx}
    &\le \int_a^b \abs{f_n(x)-f(x)}\, dx \\
    &\le \int_a^b \sup_{t\in[a,b]} \abs{f_n(t)-f(t)}\, dx \\
    &= \sup_{t\in [a,b]}\abs{f_n(t)-f(t)}\cdot \int_a^b 1\, dx
    \to 0.
  \end{align*}

  If we then define $F$ and $F_k$ as in the statement, we have
  \begin{align*}
  \Abs{F_k-F}_\infty
  &= \sup_{x\in [a,b]} \abs{\int_c^x f_k(t)-f(t)\, dt} \\
  &\le \sup_{x\in [a,b]} \abs{x-c} \cdot \Abs{f_k-f}_\infty \\
  &\le (b-a) \cdot \Abs{f_k-f}_\infty
  \to 0.
  \end{align*}
\end{proof}
  
We recall that continuous functions are locally Riemann integrable. Thus, the
previous theorem allows us to assert that the integral operator $S\colon
C^0([a,b]) \to C^0([a,b])$
\[
S(f)(x) = \int_{c}^x f(t)\, dt
\]
(which, for a fixed $c \in [a,b]$, associates to a function $f\in C^0([a,b])$
its integral function) is a continuous operator with respect to the uniform
norm.

Unfortunately, the previous theorem cannot be generally valid for improper
integrals, as we have the following counterexample.

\begin{example}\label{ex:convergenza_non_dominata}
  Consider the sequence of functions $f_n\colon \RR\to\RR$
  \[
    f_n(x) = \frac 1 n \cdot \frac{1}{1+\enclose{\frac{x-n} n}^2}.
  \]
  It can be observed that $\sup_{x\in \RR} \abs{f_n(x)} = f_n(n) = \frac 1 n$ 
  and therefore $f_n \rightrightarrows 0$ on all of $\RR$. 
  Despite this, we have 
  \[
  \int_{-\infty}^{+\infty} f_n(x) \, dx  
  = \int_{-\infty}^{+\infty} \frac{1}{1+t^2}\, dt = \pi
  \]
  and thus the integrals do not tend to zero.
\end{example}
%
In the previous example, the functions $f_n$ have a bell shape. As $n$
increases, the bell shifts towards $+\infty$, and it also lowers and widens in
such a way as to keep the integral constant. The \emph{mass} enclosed by the
bell is thus disappearing at infinity while the functions tend to zero.

To be able to pass to the limit within an improper integral, we need an
additional hypothesis: the sequence of functions must be \emph{dominated} by a
single integrable function, as we see in the following theorem. Moreover, the
hypothesis of uniform convergence can be relaxed by requiring that it holds only
on intervals strictly contained within the integration interval $(a,b)$. In this
way, the theorem also applies in the case where the interval $(a,b)$ is
unbounded or the integrand function tends to infinity at the endpoints of the
interval.

\begin{theorem}[Dominated Convergence Locally Uniform]%
  \label{th:convergenza_dominata_uniforme}%
  \mymargin{dominated uniform convergence}%
\index{dominated uniform convergence}%
  \index{theorem!of dominated convergence}%
  Let $f_k\colon (a,b) \to \RR$ and $f\colon (a,b)\to \RR$ be functions 
  locally Riemann integrable
  on the interval $(a,b)$ with $-\infty \le a < b \le +\infty$
  such that for every $[\alpha,\beta]\subset (a,b)$ 
  there is uniform convergence: $f_k\rightrightarrows f$ 
  on the interval $[\alpha,\beta]$.
  
  Suppose further that there exists $g\colon (a,b)\to \RR$
    integrable in the improper sense on $(a,b)$ with a convergent integral and
  such
  that
  \[
  \abs{f_k(x)}\le g(x)
  \]
  for every $k\in \NN$ and every $x\in (a,b)$.
    
  Then $f_k$ and $f$ have a convergent improper integral on $(a,b)$ and we have
  \[
    \lim_{k\to+\infty} \int_a^b f_k(x)\, dx = \int_a^b f(x)\, dx.
  \]
\end{theorem}
%
\begin{proof}
  Each $f_k$ is absolutely integrable in the improper sense because
  $\abs{f_k}\le g$ where $g$ has a finite integral by hypothesis.
  Since $f_k(x)\to f(x)$ for every $x\in (a,b)$
  we also have $\abs{f(x)}\le g(x)$ for every $x\in (a,b)$ 
  and thus the integral of $f$ is also absolutely convergent on $(a,b)$.
  
  Since the improper integral is convergent, we know that 
  for every $\eps>0$ there exist $\alpha>a$ and $\beta<b$ such that
  \[
    \int_\beta^b g(x)\, dx < \eps, \qquad 
    \int_a^\alpha g(x)\, dx < \eps.
  \]
    But on $[\alpha ,\beta]$ we have the convergence of the integrals by the
  previous theorem
  thus there exists $N$ such that for every $k > N$ we have 
  \[
   \abs{\int_\alpha^\beta f_k(x)\, dx - \int_\alpha^\beta f(x)\, dx} \le \eps. 
  \]
  Combining these two facts, we obtain,
  for every $k>N$
  \begin{align*}
    \MoveEqLeft{\abs{\int_a^b f_k(x)\, dx - \int_a^b f(x)\, dx}} \le \abs{\int_\alpha^\beta f_k(x)\, dx - \int_\alpha^\beta f(x)\, dx} \\
    & \quad + \abs{\int_a^\alpha (f_k(x) - f(x))\, dx }
      + \abs{\int_\beta^b (f_k(x) - f(x))\, dx } \\
    &\le \eps 
    + \int_a^\alpha \abs{f_k(x)}\, dx + \int_a^\alpha\abs{f(x)}\, dx
    + \int_\beta^b \abs{f_k(x)}\, dx + \int_\beta^b\abs{f(x)}\, dx\\
    &\le \eps 
    + 2 \int_a^\alpha g(x)\,dx
    + 2 \int_\beta^b g(x)\,dx 
    \le 5 \eps.
  \end{align*}
  We have thus verified the thesis through the definition of limit.
\end{proof}

\begin{example}
The sequence $f_k(x) = \sin(x^k)$ converges pointwise to $f(x)=0$ on the
interval
$\closeopeninterval{0}{1}$ and the convergence is uniform on every interval
$[0,\beta]$
with $\beta < 1$ (verify this!).
Moreover, $\abs{f_k(x)} = \sin(x^k) \le \sin(1)$ for every $x\in \Enclose{0,1}$
and $\int_0^1 \sin(1)\, dx = \sin(1)< +\infty$.
Thus, we can apply the dominated convergence theorem and,
by interchanging the limit with the integral, we can deduce that
\[
  \int_0^1 \sin (x^k)\, dx \to 0, \qquad \text{as $k\to +\infty$.}
\]
\end{example}

In theorem~\ref{th:convergenza_dominata}
the hypothesis that there exists a function $g$ that "dominates" the entire
sequence
is essential, as we have seen in example~\ref{ex:convergenza_non_dominata}.
Instead, the hypothesis that the sequence converges uniformly is not really
necessary:
pointwise convergence (provided the sequence is dominated) is sufficient 
to allow the passage to the limit under the integral sign.
In its generality, this result (Lebesgue's Dominated Convergence Theorem) 
is proven for the Lebesgue integral.
\index{Lebesgue!integral of}%
\index{theorem!dominated convergence}%
\index{convergence!dominated}%
\index{theorem!Lebesgue}%
\index{integral!of Lebesgue}%
However, we are able to state a slightly weaker version valid for improper
Riemann integrals that can be very useful in many practical cases.

\begin{theorem}[Dominated Convergence]
  \mymargin{dominated convergence}%
  \index{dominated convergence}%
  \label{th:convergenza_dominata}%
  \mymark{**}%
Let $I\subset \RR$ be an interval and let 
$-\infty \le a < b \le +\infty$.
Let $f\colon I \times(a,b)\to \RR$ be a continuous function 
and let $g\colon(a,b)\to \RR$ be a continuous function with a 
convergent (improper) integral
\[
 \int_a^b g(t)\, dt < +\infty
\]
such that for every $x\in I$ and for every $t\in(a,b)$, $\abs{f(x,t)} \le g(t)$.

Then for every $x\in I$, the integral $\int_a^b f(x,t)\, dt$ is convergent and
is continuous
in the variable $x$, i.e.:
\[
  \lim_{y\to x} \int_a^b f(y,t)  \, dt 
  = \int_a^b f(x,t)\, dt.
\]
\end{theorem}
%
\begin{proof}
Note that the index $k\to +\infty$, $k\in \NN$ 
in theorem~\ref{th:convergenza_dominata_uniforme} 
can be easily replaced with 
a continuous parameter $y\to x$, $x\in \RR$
thanks to theorem~\ref{th:ponte} (connection between function limit 
and sequence limit).
We therefore want to reduce to theorem~\ref{th:convergenza_dominata_uniforme} 
and to do this, we will demonstrate that for a fixed $[\alpha,\beta]\subset
(a,b)$
and a fixed $x\in I$, there is uniform convergence 
of the functions $t\mapsto f(y,t)$ to the function $t\mapsto f(x,t)$
as $y\to x$. 

Without loss of generality, we can assume that $I$ is a closed and bounded
interval
that contains the fixed point $x$. 
The function $f$ is thus continuous on the sequentially compact set $I\times
[\alpha,\beta]$
and therefore is uniformly continuous
by the Heine-Cantor Theorem~\ref{th:HeineCantor} (in two variables).
This means that for every $\eps>0$, there exists $\delta>0$ such that for every
$x,y\in I$
and for every $s,t\in[\alpha,\beta]$, we have:
\[
\sqrt{(x-y)^2+(t-s)^2} < \delta \implies \abs{f(x,t)-f(y,s)} < \eps.  
\]
But then for every $\eps>0$, if $x\in I$ and $\abs{x-y} < \delta$, we have 
\[
  \sup_{t\in\closeinterval{\alpha}{\beta}} \abs{f(y,t)-f(x,t)} \le \eps
\]
and therefore
\[
  \lim_{y\to x}  \sup_t \abs{f(y,t), f(x,t)} = 0.
\]
as we wanted to demonstrate.
\end{proof}
%
The dominated convergence theorem is often used 
to interchange the derivative with the integral as stated in the following.
%
\begin{theorem}[Derivative Passage Under the Integral Sign]%
  \mymark{**}%
  \label{th:derivata_integrale}%
  \mymargin{interchange of derivative with integral}%
\index{interchange of derivative with integral}%
Let $f\colon I\times(a,b)\to \RR$ be a continuous function, 
where $I\subset \RR$ is an interval and $-\infty \le a < b \le +\infty$.
Suppose that for every $t\in (a,b)$, the function is differentiable with respect
to $x$
for every $x\in I$
and that the derivative $\frac{\partial f}{\partial x}$ is continuous over the
entire rectangle
$I\times(a,b)$.
Suppose finally that for every $x\in I$, the integral
$\int_a^b f(x,t)\, dt$ is convergent and that
there exists a function $G\colon(a,b)\to \RR$ locally Riemann integrable 
with a convergent integral $\int_a^b G(t)\, dt$ such that for every $x\in I$ 
and for every $t\in (a,b)$, we have 
\[
  \abs{\frac{\partial f}{\partial x}(x,t)} \le G(t).
\]

Then for every $x\in I$, we have
\[
  \frac{d}{dx} \int_a^b f(x,t)\, dt = \int_a^b \frac{\partial f}{\partial x}(x,t)\, dt.  
\]
\end{theorem}
%
\begin{proof}
Fix $x_0\in I$ and define 
\[
 g(x,t) = \begin{cases}
 \frac{f(x,t) - f(x_0,t)}{x-x_0} &\text{if $x\neq x_0$}\\
 \frac{\partial f}{\partial x}(x,t) & \text{if $x = x_0$}.
\end{cases}  
\]
We then observe that the derivative of the integral at the point $x_0$ 
is:
\begin{align*}
    \Enclose{\frac{d}{dx} \int_a^b f(x,t)\, dt}_{x=x_0}
    &= \lim_{x\to x_0} \frac{\int_a^b f(x,t)\, dt - \int_a^b f(x_0,t)\, dt
  }{x-x_0}\\
  &= \lim_{x\to x_0} \int_a^b g(x,t)\, dt  
\end{align*}
while the integral of the derivative is 
$\int_a^b g(x_0,t)\, dt.$
It will suffice to demonstrate that it is possible to apply
theorem~\ref{th:convergenza_dominata}
to the function $g(x,t)$.

It is easy to verify that $G$ dominates $g$ because by the Lagrange theorem, we
know
that for every $x$ and for every $t$, there exists $x'$ with
$\abs{x'-x_0}<\abs{x-x_0}$ such that
\[
  \abs{g(x,t)} = \abs{\frac{\partial f}{\partial x}(x',t)} \le G(t).
\]

We only need to demonstrate that the function $g\colon I\times (a,b)\to \RR$ 
is continuous. 
Clearly, if $x\neq x_0$, the function is continuous; it remains to verify that
it is also continuous at the points $(x_0,t_0)$ for every $t_0\in (a,b)$.
Since by hypothesis $\partial f / \partial x$ is continuous, we know 
that 
\[
  \forall \eps>0 \colon \exists \delta>0\colon 
  {\scriptstyle\sqrt{(x-x_0)^2+(t-t_0)^2}<\delta} 
  \implies \abs{\frac{\partial f}{\partial x}(x,t) - \frac{\partial f}{\partial x}(x_0,t_0)}<\eps.
\]
Applying, as before, the Lagrange theorem, we know that if $(x,t)$ is within
$\delta$ of $(x_0,t_0)$, we have
\[
\abs{g(x,t) - g(x_0,t_0)} 
= \abs{\frac{\partial f}{\partial x}(x',t) - \frac{\partial f}{\partial
x}(x_0,t_0)}
 < \eps.
\]
This concludes the proof of the continuity of $g(x,t)$.
\end{proof}

\begin{exercise}[Derivative of the $\Gamma$ Function]
  Show that the function $\Gamma$ defined by (see example~\ref{ex:gamma_eulero})
  \[
   \Gamma(x) = \int_0^{+\infty} e^{-t}\cdot t^{x-1}\, dt 
  \]
  is differentiable and its derivative is 
  \[
  \Gamma'(x) = \int_0^{+\infty} e^{-t}\cdot \ln t \cdot t^{x-1}\, dt.  
  \]
\end{exercise}
\begin{proof}[Solution.]
It suffices to verify that we can apply theorem~\ref{th:derivata_integrale}
to the function 
\[
  f(x,t) = e^{-t}t^{x-1}
\]
whose partial derivative with respect to $x$ is 
\[
  \frac{\partial f}{\partial x}(x,t) = e^{-t}\cdot \ln t \cdot t^{x-1}.  
\]
Clearly, this function is continuous; we just need to find a function 
that dominates it. Fixing $x_0>0$, we can restrict the domain of the 
variable $x$ to the interval $\closeinterval{\frac{x_0} 2}{2x_0}$ and consider
the function 
\[
    g(t) = e^{-t} \cdot \abs{\ln t}\cdot \max\ENCLOSE{t^{\frac {x_0} 2
  -1},t^{2x_0 -1}}.
\]
This is a continuous function (as it is a composition of continuous functions, 
noting that the maximum between two continuous functions is a continuous
function).
Moreover, for $t\to 0^+$, we have 
\[
  g(t) \sim \abs{\ln t}\cdot t^{\frac{x_0}{2}-1} 
  \ll t^{-1+\eps}
\]
provided we take $\eps < \frac{x_0}{2}$.
Thus, the integral $\int_0^1 g(t)\, dt$ is convergent by 
the asymptotic comparison criterion.
If instead $t\to +\infty$, we have 
\[
  g(t) \sim e^{-t}\cdot \ln t \cdot t^{2x_0-1} \ll \frac{1}{t^2}  
\]
since the exponential tends to zero much faster 
than the logarithm and power tend to infinity.
Thus, by the asymptotic comparison criterion, 
the integral $\int_1^{+\infty} g(t)\, dt$ is also convergent.
\end{proof}

The method used in the following examples is sometimes called
 \emph{Feynman's trick}%
\mymargin{Feynman's trick}\index{trick!of Feynman}:
\index{Feynman!trick for integrals}%
\index{trick!of Feynman!for integrals}%
it involves multiplying the integrand function by a function 
that depends on a parameter so that differentiating with respect to the
parameter
simplifies the integrand function.

\begin{exercise}[Integral of the Gaussian]
  \label{ex:integrale_gaussiana}%
  \index{function!Gaussian}%
  \index{integral!Gaussian}%
  \index{Gauss!function of}%
Prove that 
\[
  \int_{-\infty}^{+\infty} e^{-s^2}\, ds = \sqrt \pi.
\]
\end{exercise}
%
\begin{proof}[Solution.]
As seen in example~\ref{ex:integrale_gaussiana_finito},
we know that the integral
\[
  I = \int_{0}^{+\infty} e^{-s^2}\, ds
\]
is finite. 
The sought integral is $2I$ (by symmetry).
Consider the following function
\[ 
  F(x) = \int_0^{+\infty} \frac{e^{-x^2(1+t^2)}}{1+t^2}\, dt.
\]
We can easily calculate 
\[
    F(0) = \int_0^{+\infty} \frac{1}{1+t^2}\, dt = \Enclose{\arctg
  t}_0^{+\infty} = \frac \pi 2.
\]
While if $x>1$, we have 
\[
  0\le \frac{e^{-x^2(1+t^2)}}{1+t^2}
  = e^{-x^2}\cdot \frac{e^{-x^2t^2}}{1+t^2}
  \le e^{-x^2}\cdot e^{-t^2}
\]
and therefore for $x\to +\infty$:
\[
  0\le F(x) \le e^{-x^2} \int_0^x e^{-t^2}\, dt \le e^{-x^2}\cdot I \to 0.
\]
Now calculate the derivative of $F$ by applying
theorem~\ref{th:derivata_integrale}.
The integrand function has a derivative 
\[
 \frac{\partial}{\partial x}\enclose {\frac{e^{-x^2(1+t^2)}}{1+t^2}}
 = -2x e^{-x^2(1+t^2)}.
\]
If we restrict the variable $x$ to an interval $\closeinterval{x_1}{x_2}$,
we have 
\[
 \abs{-2x e^{-x^2(1+t^2)}}
 \le 2 x_2 e^{-x_1^2(1+t^2)} \ll \frac{1}{1+t^2}
 \qquad\text{for } t\to+\infty
\]
and the dominating function $\frac{1}{1+t^2}$ has a finite integral over the
interval
$\openinterval{0}{+\infty}$. 
We can therefore interchange the derivative with the integral to obtain 
\[
 F'(x) = - 2 x\int_0^{+\infty} e^{-x^2(1+t^2)}\, dt
       = - 2 x e^{-x^2}\int_0^{+\infty} e^{-x^2t^2}\, dt. 
\]
With the change of variable $s=xt$, $ds=x dt$, we obtain
\[
 F'(x) = -2 e^{-x^2}\int_0^{+\infty} e^{-s^2}\, ds = - 2 e^{-x^2}\cdot I.
\]
But then on one hand, we have 
\[
  \int_0^{+\infty} F'(x)\, dx 
  = \Enclose{F(x)}_0^{+\infty}  = 0 - \frac{\pi}{2}.
\]
On the other hand, 
\[
 \int_0^{+\infty} F'(x)\, dx 
 = - 2 I \cdot \int_0^{+\infty} e^{-x^2}\, dx
 = - 2 I^2. 
\]
We conclude that $2I^2 = \frac \pi 2$, hence 
$2I=\sqrt \pi$, as we wanted to prove.
\end{proof}

\begin{exercise}[Notable Integral]
  \label{ex:integrale_sin_integrale}%
Show that 
  \[
    \int_0^{+\infty} \frac{\sin t}{t}\, dt = \frac \pi 2.
  \] 
\end{exercise}
\begin{proof}
For every $x \ge 0$, consider the function 
\[
  F(x) = \int_0^{+\infty} \frac{\sin t}{t}e^{-tx}\, dt
\]
so that the required integral is nothing but $F(0)$.
For every $x>0$,
we can differentiate the function using theorem~\ref{th:derivata_integrale}
since the derivative of the integrand function 
\[
  \frac{\partial}{\partial x}\enclose{\frac{\sin t}{t}e^{-tx}} 
  = -\sin t \cdot e^{-tx}
\]
is dominated by the function $g(t) = e^{- \eps t}$ if $x>\eps$ and 
$g$ has a convergent integral.
Thus, we have 
\[
  F'(x) = - \int_0^{+\infty} \sin t \cdot e^{-tx}\, dt. 
\]
This integral can be calculated by integrating by parts twice 
(see exercise~\ref{ex:821685}) and we find 
$F'(x) = -\frac{1}{1+x^2}$, from which
\[
  F(x) = -\arctg x + c
\]
for some $c \in \RR$.

To determine $c$, 
we now want to verify that $F(x) \to 0$ as $x\to +\infty$.
We can do this by applying Theorem~\ref{th:convergenza_dominata}, observing that
the function $f(x,t) = \frac{\sin t}{t}e^{-tx}$ is dominated by $g(t) = e^{-t}$,
for $x\ge 1$.

Thus, $c=\frac{\pi} 2$ and we have, for every $x>0$,
\[
  F(x) = \frac{\pi} 2 - \arctg x.
\]

The required integral is $F(0)$, and the formula just found tells us that
$F(x)\to \frac \pi 2$ as $x\to 0$.
It remains only to demonstrate that $F(x)\to F(0)$ as 
$x\to 0$, i.e.,
\[
  \lim_{x\to 0} \int_0^{+\infty} f(x,t) \, dt
  = \int_0^{+\infty} f(0,t) \, dt.
\]
Unfortunately, it is not possible to directly use the theorems 
of limit passage under the integral sign because 
the integral on the right side, i.e., the required integral, is not 
absolutely convergent, and therefore it will not be possible to find 
a \emph{dominating} function as required in the hypotheses of the theorem.

The idea is thus to reduce the non-absolutely convergent integral 
to an absolutely convergent integral by means of integration by parts, 
as we did in example~\ref{ex:48864}.

Indeed, we have 
\begin{align*}
F(0) - F(x) 
  = \int_0^{+\infty} \frac{\sin t}{t}\, dt  - \int_0^{+\infty} \frac{\sin
 t}{t}e^{-xt}\, dt \\
 & = \int_0^{+\infty} \frac{\sin t}{t}\enclose{1-e^{-xt}}\, dt.
\end{align*}
Now, we split the integral over the two intervals $\closeinterval{0}{\pi}$ 
and $\closeopeninterval{\pi}{+\infty}$.
On the first interval, there will be no problems: the integrand function 
indeed tends uniformly to zero, and the integral tends to zero 
by theorem~\ref{th:scambio_limite_integrale}.

On the second interval, however, we need to perform integration by parts:
\begin{align*}
  \int_\pi^{+\infty} \frac{\sin t}{t}\, dt 
  &= \Enclose{\frac{-\cos t}{t}}_{\pi}^{+\infty}
    - \int_{\pi}^{+\infty} \frac{\cos t}{t^2}\, dt \\
  &= -\frac{1}{\pi} - \int_{\pi}^{+\infty} \frac {\cos t}{t^2}\, dt.\\
\end{align*}
Then (we omit the calculation details, 
it involves integrating by parts twice as 
we did in example~\ref{ex:48864}):
\[
\int \sin t\cdot e^{-tx}\, dt 
= -h(x,t)
\]
with 
\[
h(x,t)= \frac{\cos t+x\sin t}{1+x^2}e^{-xt}.  
\]
Thus 
\begin{align*}
  \int_\pi^{+\infty} \frac{\sin t}{t} e^{-xt}\, dt 
    &= \Enclose{-\frac{\cos t+ x \sin t}{1+x^2}e^{-tx}\cdot
  \frac{1}{t}}_{\pi}^{+\infty}
  - \int_\pi^{+\infty}\frac{h(x,t)}{t^2}\, dt.\\  
    &= -\frac 1 \pi\cdot \frac{1}{1+x^2}e^{-\pi x} -
  \int_\pi^{+\infty}\frac{h(x,t)}{t^2}\, dt.
\end{align*}
and
\begin{align*}
\int_\pi^{+\infty} \frac{\sin t}{t}\, dt 
-\int_\pi^{+\infty} \frac{\sin t}{t} e^{-xt}\, dt
&= -\frac 1 \pi \enclose{1-\frac{e^{-\pi x}}{1+x^2}}
-\int_\pi^{+\infty} \frac{\cos(t) - h(x,t)}{t^2}\, dt 
\end{align*}
It is now easy to verify that if $t\in(0,+\infty)$
and $x\in[0,1]$, we have 
\[
\abs{\cos t - h(x,t)} \le \abs{\cos t} + \frac{\abs{\cos t}+x\abs{\sin t}}{1+x^2}
\le 1+ \frac{1+x}{1+x^2} \le 1 + \frac 2 1 = 3
\]
thus $\abs{\cos t - h(x,t)}/t^2 \le 3/t^2$, and having $1/t^2$ 
a convergent integral over $[\pi,+\infty)$, thanks 
to theorem~\ref{th:convergenza_dominata}, we conclude 
that 
\[
\lim_{x\to 0} \int_\pi^{+\infty} \frac{\cos t - h(x,t)}{t^2}\, dt 
= \int_{\pi}^{+\infty} \frac{\cos t - h(0,t)}{t^2}\, dt 
= 0.  
\]
But then we have demonstrated that $F(x)\to F(0)$ as $x\to 0$, 
and we have concluded the proof.
\end{proof}